Recent verifier-noise corrections already recognize that reward calibration and group normalization interact. In version~4, \citet{cai2025} explicitly shows that current-group standardization can cancel binary affine corrections, and implements direct trajectory coefficients with leave-one-out baselines and fixed or running scales. Its selective appeals mechanism uses Horvitz--Thompson counts to estimate false-negative rates, so budgeted verification within policy training is also established. The stated gradient guarantees concern an oracle class-conditional channel before finite-group normalization and clipping; the implementation additionally discusses instance-dependent residual bias and audits the advantages actually passed to training.

Version~3 of \citet{mansouri2025} likewise identifies affine invariance and corrects both reward moments and normalization. Its exact GRPO comparison requires the oracle population standard deviation; the practical corrected-variance normalizer is a consistent plug-in approximation. The paper also includes an empirical clipping ablation. These distinctions prevent interpreting our elementary counterexamples as refutations of either paper's analyzed variants.

Our estimand is the normalized or clipped update of one fixed, completely labelled finite group. We randomize reference-label queries and seek exactness conditional on that group's trajectories and potential labels, without identifying a reward-flip channel. This preserves a specified nonlinear surrogate; it does not establish that preserving that surrogate is preferable to using the simpler alternative objectives studied in prior work.
